\begin{table}[H]
\centering
\caption{Calibration-size and utility ablations for Probability-Space Retrieval Gain Calibration.}
\label{tab:ablation}
\small
\begin{adjustbox}{max width=\textwidth}
\begin{tabular}{@{}lcccccc@{}}
\toprule
Method Variant & Probability Features & Gain Calibration & Cost Penalty & F1 (\%) & Mean F1 Regret & Avg. Retrieval Calls \\
\midrule
Original Adaptive-RAG argmax & \yes & \no & \no & \NA & \NA & \NA \\
Distance-weighted KNN (smaller calibration subset) & \yes & \yes & \no & \NA & \NA & \NA \\
Distance-weighted KNN (full calibration set) & \yes & \yes & \no & \NA & \NA & \NA \\
Distance-weighted KNN + cost-aware utility & \yes & \yes & \yes & \NA & \NA & \NA \\
\bottomrule
\end{tabular}
\end{adjustbox}
\TableNote{Pending: dataset, calibration/validation/test splits, calibration-subset $N$, and held-out-test $N$ are not finalized; no result denominator is claimed. Future rows must state exact unique-question $N$ for each nested calibration subgroup and use the same held-out IDs for all metrics. \yes/\no{} mark components; tuning excludes test data, and the Oracle is not an ablation row.}
\end{table}
